% !Mode:: "TeX:UTF-8"
% 此文件从2021.10.10开始写作


%本书包括十五章以及两个附录．前两章是数学准备知识；第三章描述了多重线性代数；
%第四章到第七章分别讲述了微分流形、型式场、仿射联络以及矢量丛；
%之后详细描述了黎曼几何、测地线、子流形、超曲面、零模、复几何知识；
%继而较为详细地介绍了李群、李代数和对称空间内容；最后初步阐述了Lorentz旋量．
%
%本书需要高等数学、线性代数、数学物理方法等先修课程．
%
%本书适用于物理系高年级本科生、研究生，尤其是相对论从业者．



\section*{自序}


本书是写给理论物理工作者的，主要是为相对论学者提供一份微分几何的{\kaishu 初级}读物．
与类似的纯数学书籍相比，在深度、广度以及严谨性方面自然是逊色的，这是物理工作者撰写数学书籍的必然结果．


%依照惯例，对全书内容作如下简介．
%
%第一部分包括前三章，是微分几何必要的准备知识；若读者已熟悉这些内容，直接跳过即可．
%古典微分几何是研究嵌入到三维欧氏空间的二维正则曲面的学科，一般说来这不属于现代微分几何；
%故我们只在第一章予以简单的综述．
%之后，极扼要地介绍了微分几何所需的代数概念和点集拓扑知识．
%接着，较为详尽地给出了多重线性代数内容，并由此给出了张量的确切定义．
%
%第二部分包括四到七章．
%第四章阐述了微分流形体系；首先给出微分流形定义，并由此引入切空间、诱导映射以及子流形概念，
%其中切空间是整个微分几何随处会用到的概念；
%再者，介绍了切丛，这为后面引入仿射联络奠定了基础；
%然后，细致地描述了切矢量场和张量场的知识，这是微分流形中极为重要的概念；
%本章最后给出了单参数微分同胚群和李导数．
%第五章，首先引入外微分知识；然后简单叙述了Frobenius定理；最后给出了流形上积分的知识．
%第六章，先是定义了切丛上的仿射联络，并详尽讨论相关知识；接着定义了黎曼曲率，此曲率是整个微分几何的核心；
%然后讨论切标架丛上的黎曼曲率，这是Cartan的杰作．
%第七章（初次阅读可跳过此章），简要介绍了矢量丛（向量丛）的相关知识，
%它是切丛的推广；最后初步阐述了矢量丛与规范场的关系．
%
%第三部分包括八到十一章．
%第八章书写了黎曼几何的相关知识；首先给光滑流形配以广义黎曼度量，这样流形就成了黎曼流形；
%再者更为详尽地讨论了各种曲率（包括相对论中用到的爱因斯坦张量）；
%接着描述了等距映射和Killing切矢量场知识，这是广义相对论中所需的最基本概念；
%然后给出体积元的定义，以及Hodge星算子等概念；
%进而细致讲述了度量场下的黎曼曲率型式理论；
%最后初步讨论了截面曲率和共形映射．
%第九章阐述了黎曼流形中的测地线以及Jacobi场等相关知识，比如指数映射、法坐标系等．
%第十章先是介绍了子流形的一般内容，然后细致地讲解了超曲面知识，这是学习广义相对论所必需了解的．
%第十一章阐述了与零模（Null，类光）相关的几何，这是正定度量几何中没有的内容；
%先讲解了类光超曲面，以补上第十章所缺内容；然后详尽介绍了Newman--Penrose框架；
%最后介绍了Weyl张量的Petrov分类和Goldberg--Sachs定理．
%前十一章内容是微分几何中基础知识，是学习广义相对论必需了解的内容．
%
%第四部分包括其余章节．
%第十二章概述了复几何知识，将复几何放在此处是为李群作准备．
%第十三章是内容庞杂的一章，介绍了物理学中所需要的李群、李代数初步知识．
%第十四章讲解了对称空间内容，是广义相对论中的高阶内容．
%第十五章具体讲述了物理学中最为常用的两个李群$SU(2)$、$SL(2,\mathbb{C})$；
%进而阐述了Lorentz旋量的初步知识．



{\kaishu
    整本书中几乎没有我的原创或首创内容，但几乎每个公式我都计算过．
    唯一的原创、首创（见附录\ref{chfd-liu}）是：  
    用解存在性需求阐述了源守恒定律为何独立于场方程；
    这纠正了过去一个多世纪以来理论物理最底层存在的错误．
    具体来说，对于电磁场，电荷守恒定律独立于麦克斯韦方程组；
    麦氏散度方程（$\nabla\cdot \boldsymbol{D}=\rho,\ \nabla \cdot \boldsymbol{B}=0$）独立于麦氏旋度方程，
    它们不能看成旋度方程的初始条件．
    对于引力场，源运动方程（$\nabla^a T_{ab}=0$）独立于爱因斯坦引力场方程（$G_{ab}=8\pi T_{ab}$）．
}


严格说来，整本书是我的读书笔记．
笔者主要研读了\textcite{cc2001-zh}名著，此教材以“浓缩”著称，不易阅读．
陈维桓教授的几本教材\cite{chenwh2001,chen-li-2023-2ed-v1}相对容易多了，
本书较多参考了这几本书．
除了上述书籍，笔者还参考了\textcite{oneill1983}的专著；
此书是深入学习广义相对论的必备数学书籍．


%感谢……（待补充）
%
%
%因笔者学识所限，肯定存在这样或那样的错误；欢迎建设性的批评、指正．

\vspace{1em}
%\begin{flushright}
%刘长礼（LIUCL78@qq.com）  %%%%, ORCID: 0000-0002-3592-2449 ）
%
%202x年x月于北京
%\end{flushright}
%%%\vspace{1em}






%\newpage
%
%\section*{指标惯例}
%本书主要采用由Penrose创立的抽象指标记号；同时也采用数学家熟悉的记号；但在同一章内只用一种记号．
%角标一般是$1\leqslant i, \mu \leqslant m$．
%
%
%\section*{部分数学符号表}
%\begin{tabular}{ll}
%    $\mathbb{R},\, \mathbb{C}$ & 全体实数，实数域；全体复数，复数域 \\
%    $\bar{z},\,z^*$ & 数$z$的复共轭 \\
%    $\overset{def}{=}$ & 定义，有时候“$\equiv$”也有定义的含义 \\
%    $f:A\to B$、$A\xrightarrow{f}B$ & 定义从$A$到$B$的映射$f$ \\
%    $\cong$ & 同构 \\
%    $\otimes$ & 张量积 \\
%    $\circ$ & 复合映射．例如$f\circ g$代表先作用$g$再作用$f$ \\
%    $C^r_p(U)$ & $p\in U$点的全体$r$阶连续可微函数集合 \\
%    $V^*$ & 矢量空间$V$的对偶空间 \\
%    $\frac{\partial }{\partial x^\mu}$或$(\frac{\partial }{\partial x^\mu})^a$ &  第$\mu$个坐标基矢 \\
%    ${\rm d}x^\mu$或$({\rm d}x^\mu)^a$ & 第$\mu$个对偶坐标基矢 \\
%    $\mathfrak{T}^p_q(V)$ & 矢量空间$V$上全体$(p,q)$型张量场 \\
%    $\mathfrak{X}(M),\, \mathfrak{X}^*(M)$ & 流形$M$上的全体切矢量场，余切矢量场 \\
%    $\nabla,\, {\rm D}$ & 联络 \\
%    $\phi^*$ & 映射$\phi$诱导出的余切映射或拉回映射 \\
%    $\phi_*$ & 映射$\phi$诱导出的切映射或推前映射 \\
%    $\ln, \log$ & 以$e$为底的对数 \\
%    ${\rm Span},\, {\rm Span}_{\mathbb{F}}$ & 组合系数是$\mathbb{F}$中的常数 \\
%    ${\rm Span}_{\infty}$ & 组合系数是$C^{\infty}(M)$中的函数 \\
%\end{tabular}





\endinput
